\BeleskeID{02-s031}
% element: 02-s031:e01
% element: 02-s031:e02
% element: 02-s031:e03
% element: 02-s031:d01

Prvo formiramo proizvod \(CB\), pa ga sa \(A\) dovodimo na završno NI kolo. Kada koristimo samo NI kola, proizvod se dobija u dva koraka: \(u=\overline{CB}\), a zatim \(v=\overline{uu}=CB\). Završni izlaz je \(F=\overline{vA}=\overline{CBA}\). Drugo kolo je NI upotrebljeno kao invertor; bez njega unutrašnja negacija ostala bi u izrazu.

Isti princip važi za NILI realizaciju: \(u=\overline{C+B}\), \(v=\overline{u+u}=C+B\), pa je \(F=\overline{v+A}\). Invertovanje i pravilno sačuvan međurezultat omogućavaju da se višeulazne funkcije sastave iz dvoulaznih kapija iste vrste. Možemo koristiti i samo troulazna NI ili samo troulazna NILI kola: ponavljanjem ulaza dobijamo dvoulaznu funkciju, a vezivanjem sva tri ulaza na isti signal dobijamo invertor. Pri takvom vezivanju proveravaju se i električni zahtevi.
